% =====================================================================
%  TEMPLATE MAKALAH SMA — LaTeX
%  Cover: Judul di atas logo + footer tahun ajaran
%  Compile: pdflatex makalah.tex   (jalankan 2x untuk daftar isi)
% =====================================================================
\documentclass[12pt,a4paper]{article}

% ---------- Paket ----------
\usepackage[utf8]{inputenc}
\usepackage[T1]{fontenc}
\usepackage[indonesian]{babel}
\usepackage[a4paper,top=3cm,bottom=3cm,left=4cm,right=3cm]{geometry}
\usepackage{graphicx}
\usepackage{setspace}
\usepackage{titlesec}
\usepackage{enumitem}
\usepackage{hyperref}
\usepackage{parskip}

% ---------- Pengaturan umum ----------
\onehalfspacing                      % spasi 1,5
\setlength{\parindent}{1.27cm}       % indentasi paragraf
\graphicspath{{assets/}}

\hypersetup{
  colorlinks=true,
  linkcolor=black,
  urlcolor=blue,
  citecolor=black,
  pdftitle={Makalah},
}

% Judul bab/subbab ala makalah SMA
\titleformat{\section}{\normalfont\large\bfseries}{\thesection.}{0.5em}{}
\titleformat{\subsection}{\normalfont\normalsize\bfseries}{\thesubsection}{0.5em}{}

% =====================================================================
%  GANTI DATA DI BAWAH INI
% =====================================================================
\newcommand{\judulMakalah}{PENGARUH MEDIA SOSIAL TERHADAP PRESTASI BELAJAR SISWA}
\newcommand{\jenisKarya}{Makalah}
\newcommand{\mataPelajaran}{Mata Pelajaran: Biologi}
\newcommand{\namaSiswa}{Nama Lengkap Siswa}
\newcommand{\kelasSiswa}{Kelas XI IPA 1}
\newcommand{\nomorAbsen}{No. Absen: 01}
\newcommand{\namaSekolah}{SMA NEGERI 1 CONTOH}
\newcommand{\namaGuru}{Nama Guru Pembimbing, S.Pd.}
\newcommand{\tahunAjaran}{2024/2025}
\newcommand{\namaKota}{Jakarta}

% Logo: taruh file di assets/logo.png
\newcommand{\logoFile}{assets/logo.png}

% =====================================================================
\begin{document}

% ---------------------- HALAMAN COVER ----------------------
\begin{titlepage}
  \centering
  \singlespacing

  % Judul di paling atas
  {\large\bfseries \jenisKarya\par}
  \vspace{0.8cm}
  {\Large\bfseries \MakeUppercase{\judulMakalah}\par}
  \vspace{1.2cm}

  % Logo
  \IfFileExists{\logoFile}
    {\includegraphics[width=4.5cm]{\logoFile}}
    {\fbox{\parbox[c][4.5cm][c]{4.5cm}{\centering\itshape\small Logo Sekolah\\\texttt{assets/logo.png}}}}
  \vspace{1.2cm}

  % Disusun oleh
  {\large Disusun oleh:\par}
  \vspace{0.4cm}
  {\large\bfseries \namaSiswa\par}
  {\large \kelasSiswa\par}
  {\large \nomorAbsen\par}
  \vspace{1.2cm}

  % Instansi
  {\large\bfseries \MakeUppercase{\namaSekolah}\par}
  \vspace{0.3cm}
  {\large \mataPelajaran\par}
  \vspace{0.3cm}
  {\large Guru Pembimbing: \namaGuru\par}

  \vfill

  % Footer tahun ajaran di paling bawah
  {\large\bfseries TAHUN AJARAN \tahunAjaran\par}
  \vspace{0.5cm}
\end{titlepage}

% ---------------------- KATA PENGANTAR ----------------------
\section*{Kata Pengantar}
\addcontentsline{toc}{section}{Kata Pengantar}
Puji syukur kami panjatkan kepada Tuhan Yang Maha Esa atas selesainya penyusunan \jenisKarya\ yang berjudul \textit{\judulMakalah}. \jenisKarya\ ini disusun sebagai salah satu tugas \mataPelajaran.

Kami mengucapkan terima kasih kepada \namaGuru\ selaku guru pembimbing, serta semua pihak yang telah membantu penyusunan makalah ini. Kami menyadari makalah ini masih jauh dari sempurna, sehingga kritik dan saran yang membangun sangat kami harapkan.

\vspace{1cm}
\begin{flushright}
\namaKota, \today\\[1.5cm]
\namaSiswa
\end{flushright}

\newpage

% ---------------------- DAFTAR ISI ----------------------
\renewcommand{\contentsname}{Daftar Isi}
\tableofcontents
\newpage

% ---------------------- BAB I ----------------------
\section{Pendahuluan}
\subsection{Latar Belakang}
Tuliskan latar belakang masalah di sini. Jelaskan fenomena, data, atau fakta yang melandasi pemilihan topik makalah.

\subsection{Rumusan Masalah}
\begin{enumerate}[label=\arabic*., leftmargin=*]
  \item Apa rumusan masalah pertama?
  \item Apa rumusan masalah kedua?
\end{enumerate}

\subsection{Tujuan Penulisan}
\begin{enumerate}[label=\arabic*., leftmargin=*]
  \item Menjelaskan tujuan pertama.
  \item Menjelaskan tujuan kedua.
\end{enumerate}

% ---------------------- BAB II ----------------------
\section{Pembahasan}
\subsection{Kajian Teori}
Uraikan teori-teori yang mendukung topik makalah.

\subsection{Analisis}
Uraikan analisis atau hasil pembahasan berdasarkan teori pada subbab sebelumnya.

% ---------------------- BAB III ----------------------
\section{Penutup}
\subsection{Kesimpulan}
\begin{enumerate}[label=\arabic*., leftmargin=*]
  \item Kesimpulan utama dari pembahasan.
  \item Kesimpulan pendukung.
\end{enumerate}

\subsection{Saran}
Tuliskan saran untuk pembaca, penulis, maupun pihak terkait.

% ---------------------- DAFTAR PUSTAKA ----------------------
\newpage
\section*{Daftar Pustaka}
\addcontentsline{toc}{section}{Daftar Pustaka}
\begin{enumerate}[label={}, leftmargin=*]
  \item Nama Penulis. \textit{Judul Buku}. Kota: Penerbit, Tahun.
  \item Nama Penulis. ``Judul Artikel.'' \textit{Nama Situs}, Tanggal. URL.
\end{enumerate}

\end{document}
